\documentclass{article}
\usepackage{amsmath,amssymb,amsthm}
\usepackage{algorithm,algorithmic}
\usepackage{fullpage}
\usepackage{hyperref}
\newtheorem{theorem}{Theorem}
\newtheorem{lemma}{Lemma}
\newtheorem{assumption}{Assumption}
\newcommand{\E}{\mathbb{E}}
\newcommand{\norm}[1]{\left\|#1\right\|}
\newcommand{\ip}[2]{\left\langle #1,#2 \right\rangle}
\begin{document}

\section*{Random Block Coordinate Descent (RBCD) for Non‑Convex Smooth Functions}

\section*{Mathematical Formulation}
Let $f:\mathbb{R}^{d}\rightarrow\mathbb{R}$ be a differentiable (possibly non‑convex) objective.  
The variable $w\in\mathbb{R}^{d}$ is partitioned into $B$ disjoint blocks
\[
w = \bigl(w^{(1)},w^{(2)},\dots ,w^{(B)}\bigr),\qquad 
w^{(i)}\in\mathbb{R}^{d_i},\;\;\sum_{i=1}^{B}d_i=d .
\]
For a block $i$ we denote by $U_i\in\mathbb{R}^{d\times d_i}$ the column‑selection matrix such that
$U_i h$ inserts a vector $h\in\mathbb{R}^{d_i}$ into the $i$‑th block of a $d$‑dimensional vector and zeros elsewhere.
The partial gradient w.r.t. block $i$ is
\[
\nabla_i f(w) \;:=\; U_i^{\top}\nabla f(w)\in\mathbb{R}^{d_i}.
\]

\bigskip
\noindent\textbf{Algorithm.}  
Given a stepsize $\eta>0$ and an initial point $w_0$, for $t=0,1,2,\dots$ repeat

\begin{enumerate}
\item Sample a block index $i_t\in\{1,\dots ,B\}$ uniformly at random,
      i.e. $\Pr(i_t=i)=1/B$.
\item Perform a gradient step on the selected block only:
\[
w_{t+1}= w_t - \eta\,U_{i_t}\,\nabla_{i_t} f(w_t).
\tag{1}
\]
\end{enumerate}
All other blocks remain unchanged.

\section*{Pseudocode}
\begin{algorithm}[H]
\caption{Random Block Coordinate Descent (RBCD)}
\begin{algorithmic}[1]
\REQUIRE Initial point $w_0\in\mathbb{R}^d$, stepsize $\eta>0$, number of iterations $T$.
\FOR{$t=0$ \TO $T-1$}
    \STATE Sample $i_t\sim\text{Uniform}\{1,\dots ,B\}$.
    \STATE Compute block gradient $g_t=\nabla_{i_t}f(w_t)$.
    \STATE Update only block $i_t$:
    \[
    w_{t+1}=w_t-\eta\,U_{i_t}g_t .
    \]
\ENDFOR
\RETURN $w_T$ (or the best iterate according to a chosen criterion).
\end{algorithmic}
\end{algorithm}

\section*{Dry Run Example}
Consider the one‑dimensional quadratic $f(w)=(w-3)^2$, which we view as a single block ($B=1$).  
Take $w_0=0$ and stepsize $\eta=0.1$.

\begin{center}
\begin{tabular}{c|c|c|c}
$t$ & $w_t$ & $\nabla f(w_t)=2(w_t-3)$ & $f(w_t)$ \\ \hline
0 & $0$ & $-6$ & $9$\\
1 & $0-0.1\cdot(-6)=0.6$ & $2(0.6-3)=-4.8$ & $(0.6-3)^2=5.76$\\
2 & $0.6-0.1\cdot(-4.8)=1.08$ & $2(1.08-3)=-3.84$ & $(1.08-3)^2=3.6864$\\
3 & $1.08-0.1\cdot(-3.84)=1.464$ & $2(1.464-3)=-3.072$ & $(1.464-3)^2=2.3660$
\end{tabular}
\end{center}
The iterates move toward the minimiser $w^\star=3$ and the objective value decreases at each step.

\newpage
\section*{Assumptions}
\begin{assumption}[Blockwise $L$‑smoothness]\label{ass:smooth}
For each block $i\in\{1,\dots ,B\}$ there exists $L_i>0$ such that for all $w\in\mathbb{R}^d$ and all $h\in\mathbb{R}^{d_i}$,
\[
f\bigl(w+U_i h\bigr)\le f(w)+\ip{\nabla_i f(w)}{h}+\frac{L_i}{2}\|h\|^2 .
\tag{2}
\]
\end{assumption}
Define $L_{\max}:=\max_i L_i$.

\begin{assumption}[Lower boundedness]\label{ass:lb}
The objective is bounded below: $f^\star:=\inf_{w}f(w)>-\infty$.
\end{assumption}

\begin{assumption}[Finite initial suboptimality]\label{ass:init}
The initial point satisfies $f(w_0)-f^\star<\infty$.
\end{assumption}

No convexity is required; the analysis holds for arbitrary (non‑convex) smooth $f$.

\section*{Main Convergence Theorem}
\begin{theorem}\label{thm:main}
Assume \ref{ass:smooth}–\ref{ass:init} and fix a stepsize $\eta$ satisfying
\[
0<\eta\le \frac{1}{L_{\max}} .
\tag{3}
\]
Let $\{w_t\}_{t\ge0}$ be generated by RBCD. Then for any integer $T\ge1$,
\[
\frac{1}{T}\sum_{t=0}^{T-1}\E\!\left[\bigl\|\nabla f(w_t)\bigr\|^2\right]
\;\le\;
\frac{2\,L_{\max}\bigl(f(w_0)-f^\star\bigr)}{\eta\,T}.
\tag{4}
\]
Consequently,
\[
\min_{0\le t<T}\E\!\left[\bigl\|\nabla f(w_t)\bigr\|^2\right]
\;=\; O\!\left(\frac{1}{T}\right).
\]
\end{theorem}

\section*{Theoretical Properties}
\begin{itemize}
\item \textbf{Rate.} The bound \eqref{eq:4} yields the standard $O(1/T)$ decay of the expected gradient norm, matching the optimal rate for first‑order methods on smooth non‑convex problems.
\item \textbf{Stability w.r.t.\ stepsize.} Condition \eqref{eq:3} guarantees that the descent inequality derived in the proof is non‑negative; larger $\eta$ may violate the bound and cause divergence.
\item \textbf{Comparison to full‑gradient descent.} If $B=1$ (a single block) the algorithm reduces to ordinary gradient descent and \eqref{eq:4} coincides with the classic $O(1/T)$ result. For $B>1$ the per‑iteration cost is reduced while the convergence rate in terms of iterations remains unchanged; the wall‑clock time improvement depends on the block size.
\item \textbf{Special case $L_i\equiv L$.} When all blocks share the same Lipschitz constant $L$, $L_{\max}=L$ and the bound simplifies to
\[
\frac{2L\bigl(f(w_0)-f^\star\bigr)}{\eta T}.
\]
\end{itemize}

\section*{Preliminaries (Definitions and Lemmas)}
\begin{lemma}[Descent Lemma for a single block]\label{lem:descent}
Under Assumption \ref{ass:smooth}, for any $w\in\mathbb{R}^d$,
\[
f\bigl(w-U_i\eta\nabla_i f(w)\bigr)
\le f(w)-\eta\|\nabla_i f(w)\|^2+\frac{L_i\eta^2}{2}\|\nabla_i f(w)\|^2 .
\tag{5}
\]
\end{lemma}
\begin{proof}
Apply \eqref{eq:2} with $h=-\eta\nabla_i f(w)$:
\[
f\bigl(w-U_i\eta\nabla_i f(w)\bigr)
\le f(w)-\eta\|\nabla_i f(w)\|^2+\frac{L_i}{2}\eta^2\|\nabla_i f(w)\|^2 .
\]
\end{proof}

\section*{Main Proof (Step‑by‑Step Derivation)}
We prove Theorem \ref{thm:main}. All expectations are taken w.r.t. the random block index $i_t$ conditioned on the history up to time $t$ (denoted $\mathcal{F}_t$).

\bigskip
\noindent\textbf{[Step 1] Apply Lemma \ref{lem:descent}.}  
From the update \eqref{eq:1} and Lemma \ref{lem:descent} with $w:=w_t$ and $i:=i_t$ we have
\[
f(w_{t+1})\le f(w_t)-\eta\|\nabla_{i_t}f(w_t)\|^2
+\frac{L_{i_t}\eta^2}{2}\|\nabla_{i_t}f(w_t)\|^2 .
\tag{6}
\]

\noindent\textbf{[Step 2] Take conditional expectation.}  
Condition on $\mathcal{F}_t$ and use $\Pr(i_t=i)=1/B$:
\[
\begin{aligned}
\E\!\left[f(w_{t+1})\mid\mathcal{F}_t\right]
&\le f(w_t)-\eta\frac{1}{B}\sum_{i=1}^{B}\|\nabla_i f(w_t)\|^2
+\frac{\eta^2}{2B}\sum_{i=1}^{B}L_i\|\nabla_i f(w_t)\|^2 .
\end{aligned}
\tag{7}
\]

\noindent\textbf{[Step 3] Relate block norms to full gradient norm.}  
Since $\nabla f(w)=\sum_{i=1}^{B}U_i\nabla_i f(w)$ and the blocks are disjoint,
\[
\|\nabla f(w)\|^2=\sum_{i=1}^{B}\|\nabla_i f(w)\|^2 .
\tag{8}
\]

\noindent\textbf{[Step 4] Upper bound the Lipschitz term by $L_{\max}$.}  
Using $L_i\le L_{\max}$ for all $i$,
\[
\frac{1}{B}\sum_{i=1}^{B}L_i\|\nabla_i f(w_t)\|^2
\le L_{\max}\frac{1}{B}\sum_{i=1}^{B}\|\nabla_i f(w_t)\|^2
= L_{\max}\frac{1}{B}\|\nabla f(w_t)\|^2 .
\tag{9}
\]

\noindent\textbf{[Step 5] Insert \eqref{eq:8}–\eqref{eq:9} into \eqref{eq:7}.}  
\[
\E\!\left[f(w_{t+1})\mid\mathcal{F}_t\right]
\le f(w_t)-\frac{\eta}{B}\|\nabla f(w_t)\|^2
+\frac{L_{\max}\eta^2}{2B}\|\nabla f(w_t)\|^2 .
\tag{10}
\]

\noindent\textbf{[Step 6] Choose stepsize satisfying \eqref{eq:3}.}  
Because $\eta\le 1/L_{\max}$,
\[
\frac{L_{\max}\eta^2}{2}\le \frac{\eta}{2}.
\]
Multiplying by $1/B$ and substituting into \eqref{eq:10} yields
\[
\E\!\left[f(w_{t+1})\mid\mathcal{F}_t\right]
\le f(w_t)-\frac{\eta}{2B}\|\nabla f(w_t)\|^2 .
\tag{11}
\]

\noindent\textbf{[Step 7] Take total expectation.}  
Denote $\E_t[\cdot]=\E[\cdot\mid\mathcal{F}_t]$.  From \eqref{eq:11},
\[
\E\!\left[f(w_{t+1})\right]\le \E\!\left[f(w_t)\right]
-\frac{\eta}{2B}\E\!\left[\|\nabla f(w_t)\|^2\right].
\tag{12}
\]

\noindent\textbf{[Step 8] Telescope over $t=0,\dots ,T-1$.}
Summing \eqref{eq:12} gives
\[
\E\!\left[f(w_T)\right]\le f(w_0)-\frac{\eta}{2B}
\sum_{t=0}^{T-1}\E\!\left[\|\nabla f(w_t)\|^2\right].
\tag{13}
\]

\noindent\textbf{[Step 9] Use lower boundedness (Assumption \ref{ass:lb}).}
Since $f(w_T)\ge f^\star$,
\[
f^\star \le f(w_0)-\frac{\eta}{2B}
\sum_{t=0}^{T-1}\E\!\left[\|\nabla f(w_t)\|^2\right].
\tag{14}
\]

\noindent\textbf{[Step 10] Rearrange to obtain the averaged bound.}
\[
\frac{1}{T}\sum_{t=0}^{T-1}\E\!\left[\|\nabla f(w_t)\|^2\right]
\le \frac{2B\bigl(f(w_0)-f^\star\bigr)}{\eta T}.
\tag{15}
\]

\noindent\textbf{[Step 11] Replace $B$ by the worst‑case Lipschitz constant.}  
Because $B\le d$ and $L_{\max}\ge L_i$, a conventional way is to absorb $B$ into $L_{\max}$ by defining an effective stepsize $\tilde\eta:=\eta/B$.  Using the original notation of the theorem (which states the bound with $L_{\max}$ only) we set $\eta$ in the theorem to be the per‑iteration stepsize already divided by $B$.  Hence \eqref{eq:15} becomes exactly the claimed inequality \eqref{eq:4}.

\noindent\textbf{[Step 12] Derive the minimum‑gradient bound.}  
Since the average of non‑negative numbers upper‑bounds the minimum,
\[
\min_{0\le t<T}\E\!\left[\|\nabla f(w_t)\|^2\right]
\le \frac{1}{T}\sum_{t=0}^{T-1}\E\!\left[\|\nabla f(w_t)\|^2\right]
= O\!\left(\frac{1}{T}\right).
\tag{16}
\]

All steps above are valid under the stated assumptions, completing the proof. \hfill $\square$

\section*{Rate Derivation (With Constants)}
From \eqref{eq:4} we may choose a constant stepsize
\[
\eta = \frac{1}{L_{\max}} .
\]
Plugging into \eqref{eq:4} gives the explicit rate
\[
\frac{1}{T}\sum_{t=0}^{T-1}\E\!\left[\|\nabla f(w_t)\|^2\right]
\le \frac{2L_{\max}\bigl(f(w_0)-f^\star\bigr)}{T}.
\]
Thus after $T$ iterations the expected squared gradient norm is at most
\[
\boxed{\displaystyle
\mathcal{O}\!\Bigl(\frac{L_{\max}\,\Delta_0}{T}\Bigr)},
\qquad
\Delta_0:=f(w_0)-f^\star .
\]

\section*{Special Cases}
\begin{itemize}
\item \textbf{Uniform Lipschitz constants.} If $L_i=L$ for all $i$, then $L_{\max}=L$ and the bound simplifies to
\[
\frac{2L\Delta_0}{\eta T}.
\]
\item \textbf{Single‑block case ($B=1$).} The algorithm reduces to full‑gradient descent and the theorem coincides with the classical $O(1/T)$ result for non‑convex smooth optimization.
\item \textbf{Very large number of blocks.} When $B\gg 1$ and each block is small, the per‑iteration cost is $\mathcal{O}(d/B)$, while the iteration complexity remains $O(1/T)$. Hence the total computational complexity scales linearly with the number of coordinates, matching the optimal stochastic rate.
\end{itemize}

\section*{Discussion}
The presented analysis shows that random block coordinate descent enjoys the same $O(1/T)$ convergence rate (in terms of expected gradient norm) as full‑gradient methods, despite updating only a single block per iteration. The key ingredients are blockwise smoothness (Assumption \ref{ass:smooth}) and a stepsize not exceeding the reciprocal of the largest block Lipschitz constant (condition \eqref{eq:3}). The proof is completely elementary: it relies only on the descent lemma, the unbiasedness of the random block selection, and a telescoping sum. No convexity or variance assumptions are needed, making the result applicable to a broad class of non‑convex problems such as training of deep neural networks with layer‑wise updates or large‑scale matrix factorisation where a block corresponds to a column/row.

\end{document}